\documentclass[11pt,a4paper]{article}
\usepackage[T1]{fontenc}
\usepackage{lmodern}
\catcode`\_=12
\usepackage{listings}

\usepackage[utf8]{inputenc}
\usepackage{amsmath}
\usepackage{amssymb}
\DeclareUnicodeCharacter{211D}{\ensuremath{\mathbb{R}}}
\DeclareUnicodeCharacter{2022}{\ensuremath{\bullet}}

\usepackage{amsfonts}
\usepackage{graphicx}
\usepackage{hyperref}

\title{\textbf{TONE: A Physical Core Based on Jerk-Tensor Dynamics within Discrete Correlation Matrices}}
\author{\textbf{Johannes Geisen}}
\date{October 2026}

\begin{document}

\maketitle

\begin{abstract}
This paper presents the foundational mathematical framework of TONE (Theory of Everything), demonstrating how space-time geometry and gravitational dynamics emerge from a fundamentally discrete substratum. By applying a discrete Jerk-Tensor formulation directly to emergent correlation matrices, we expand upon the principles of background-independent \textit{Internal Relativity}.
\end{abstract}
\clearpage
\section{Introduction \& Background}
In conventional general relativity, the smooth manifold of space-time is taken as a background. Conversely, approaches like Dreyer's \textit{Internal Relativity} suggest that geometry emerges from the quantum correlations of a discrete system. TONE unifies these views by introducing a kinetic drive: the \textbf{Jerk-Tensor}.

\section{Mathematical Formulation}
Let $G_t$ be a discrete quantum-like substrate evolving at discrete time steps $t$. The emergent space-time relations are captured by the correlation matrix $C_t$:
\begin{equation}
C_t = \text{corr}(G_t)
\end{equation}

While velocity and acceleration describe first and second-order changes, the structural evolution of emergent geometry requires a third-order temporal derivative. 

\subsection{The Discrete Jerk-Tensor}
Because the fundamental reality is discrete, we replace continuous time derivatives with finite differences. The discrete \textbf{Jerk-Tensor} $J_t$ acting on the correlation space is defined as the third discrete difference:
\begin{equation}
J_t = \Delta^3 C_t = C_t - 3C_{t-1} + 3C_{t-2} - C_{t-3}
\end{equation}

Where:
\begin{itemize}
    \item $C_t$ is the present emergent metric state.
    \item $J_t$ governs the acceleration-change of quantum correlations, eliminating the need for a fixed classical background.
\end{itemize}
\clearpage
\section{Computational Implementation}
The validity of this framework is demonstrated via the Python implementation in this repository (\texttt{toe\_simulation.py}), which tracks the multi-step history of the discrete substratum to dynamically compute $J\_t$.

\noindent\rule{\textwidth}{0.4pt}
\begin{center}
\textbf{TONE CORE: FORMAL MATHEMATICAL PROOF STRUCTURE FOR THE DISCRETE JERK TENSOR} \\
Intellectual Property of Johannes Geisen (frintroper)
\end{center}
\noindent\rule{\textwidth}{0.4pt}

\begin{lstlisting}[basicstyle=\small\ttfamily, breaklines=true, frame=single, xleftmargin=10pt, xrightmargin=10pt]
import Mathlib.Data.Matrix.Basic
import Mathlib.Data.Real.Basic

-- Definiere den mathematischen Raum fuer unsere diskreten Korrelationsmatrizen
def CorrelationMatrix (n : Type) [Fintype n] [DecidableEq n] := Matrix n n \mathbb{R}

-- Die mathematische Definition des diskreten Jerk-Tensors (3. Differenz) in Lean
def discrete_jerk_tensor {n : Type} [Fintype n] [DecidableEq n]
  (C_t : CorrelationMatrix n)
  (C_tm1 : CorrelationMatrix n)
  (C_tm2 : CorrelationMatrix n)
  (C_tm3 : CorrelationMatrix n) : CorrelationMatrix n :=
  C_t - (3 * C_tm1) + (3 * C_tm2) - C_tm3

-- Ein formaler Check: Wenn sich das System im Gleichgewicht befindet
-- (konstante Matrix), muss der Jerk-Tensor mathematisch exakt Null ergeben.
theorem jerk_of_constant_system_is_zero {n : Type} [Fintype n] [DecidableEq n]
  (C : CorrelationMatrix n) :
  discrete_jerk_tensor C C C C = 0 := by
  simp [discrete_jerk_tensor]
  ring
\end{lstlisting}

\clearpage

\section{Intellectual Property \& License Agreement}
\textbf{Copyright (c) 2026 Johannes Geisen (https://github.com/frintroper). All Rights Reserved.}

\subsection{Ownership of the Framework}
The mathematical logic, algebraic derivations, algorithmic architectures, and computational implementations regarding the TONE (Theory of Everything) Physical Core --- specifically the application of a third-order discrete temporal difference (Jerk-Tensor) onto emergent correlation spaces --- represent the exclusive intellectual property of Johannes Geisen.

\subsection{Strict Terms and Restrictions}
\begin{itemize}
    \item \textbf{Prohibition of Unauthorized Replication:} Copying, modifying, distributing, or sub-licensing this framework, its source code, or its underlying mathematical formulations for commercial or non-commercial deployment is strictly prohibited without prior written consent from the author.
    \item \textbf{Cross-Language Porting Restriction:} This system is legally protected against unauthorized translation, rewrites, or cross-compilation into alternative programming languages, including but not limited to C++, Rust, Lean, Julia, or Go. Any implementation outside the original repository requires explicit licensing.
    \item \textbf{Automated Scraping and AI Training Ban:} Automated web scrapers, data-mining agents, and Large Language Models (LLMs) are strictly forbidden from utilizing the contents of this repository for training purposes, model refinement, or automated derivative code generation.
\end{itemize}
\clearpage
\section{Formal Verification via Interactive Theorem Proving}
To establish mathematical infallibility and protect the TONE core against logical discrepancies, the discrete dynamics are formally verified using the Lean 4 interactive theorem prover. 

\subsection{Lean 4 Implementation Code}
The following formal specification defines the discrete Jerk-Tensor over a finite correlation matrix space and proves that a static system yields a vanishing kinetic drive:

\begin{lstlisting}[basicstyle=\small\ttfamily, breaklines=true, frame=single]
-- =====================================================
-- TONE CORE: FORMAL MATHEMATICAL PROOF STRUCTURE FOR THE DISCRETE JERK TENSOR
-- Intellectual Property of Johannes Geisen (frintroper)
-- =====================================================

import Mathlib.Data.Matrix.Basic
import Mathlib.Data.Real.Basic

-- Define the mathematical space for discrete correlation matrices
def CorrelationMatrix (n : Type) [Fintype n] [DecidableEq n] := Matrix n n \mathbb{R}


-- Formal definition of the discrete Jerk-Tensor (3rd difference)
def discrete_jerk_tensor {n : Type} [Fintype n] [DecidableEq n] 
  (C_t : CorrelationMatrix n) 
  (C_tm1 : CorrelationMatrix n) 
  (C_tm2 : CorrelationMatrix n) 
  (C_tm3 : CorrelationMatrix n) : CorrelationMatrix n :=
    C_t - (3 * C_tm1) + (3 * C_tm2) - C_tm3


-- Proof: If the system is in perfect equilibrium, the Jerk-Tensor evaluates to zero
theorem jerk_of_constant_system_is_zero {n : Type} [Fintype n] [DecidableEq n] 
  (C : CorrelationMatrix n) : 
  discrete_jerk_tensor C C C C = 0 := by
  simp [discrete_jerk_tensor]
  ring
\end{lstlisting}


\end{document}
